\documentclass[11pt,a4paper]{article}
\usepackage[utf8]{inputenc}
\usepackage[T1]{fontenc}
\usepackage{lmodern}
\renewcommand{\familydefault}{\sfdefault}
\usepackage[a4paper,margin=2cm,top=2.2cm,bottom=2.2cm]{geometry}
\usepackage{xcolor}
\usepackage{tabularx}
\usepackage{booktabs}
\usepackage{array}
\usepackage{enumitem}
\usepackage{titlesec}
\usepackage{fancyhdr}
\usepackage[hidelinks]{hyperref}

\definecolor{brand}{HTML}{1F3A5F}
\definecolor{panel}{HTML}{EEF3F8}
\definecolor{muted}{HTML}{5A6570}

\titleformat{\section}{\large\bfseries\color{brand}}{\thesection.}{0.5em}{}
\titlespacing*{\section}{0pt}{1.1em}{0.4em}
\setlist{nosep,leftmargin=1.4em}
\setlength{\parindent}{0pt}
\setlength{\parskip}{0.45em}
\renewcommand{\arraystretch}{1.25}

\pagestyle{fancy}
\fancyhf{}
\renewcommand{\headrulewidth}{0pt}
\fancyfoot[L]{\footnotesize\color{muted}NAPQES --- Independent Security Audit Summary}
\fancyfoot[R]{\footnotesize\color{muted}Epineon $\cdot$ October 2026 $\cdot$ \thepage}

\newcommand{\pquote}[1]{\par\smallskip{\color{brand}\leftskip=1.2em\rightskip=1.2em\itshape ``#1''\par}\smallskip}

\begin{document}

{\color{brand}\Huge\bfseries NAPQES}\\[0.2em]
{\color{brand}\LARGE Independent Security Audit Summary}\\[0.4em]
{\color{muted}For stakeholders, distributors and customers \quad$\cdot$\quad Epineon \quad$\cdot$\quad October 2026}

\medskip
\noindent\colorbox{panel}{\begin{minipage}{\dimexpr\linewidth-2\fboxsep}
\smallskip
\textbf{\color{brand}In short.} NAPQES, Epineon's encryption scheme, was independently reviewed by
\textbf{ABDK Consulting} in \textbf{four review rounds} between July and October 2026. The rounds
covered both the published specification and the reference software. The findings dropped in
severity from round to round: the final round reported \textbf{no critical and no major issues}.
The auditors concluded that \emph{``the v8 encryption path is sound, and we found no way to break
it from outside.''} The protocol described in the specification is the one Epineon implements
and ships.
\smallskip
\end{minipage}}

\section{What NAPQES is}
NAPQES (\emph{Noise-Augmented Post-Quantum Encryption Scheme}) protects data in two ways at once:
\begin{itemize}
  \item \textbf{Confidentiality.} Only the holder of the key can read the data.
  \item \textbf{Integrity.} Any change to the protected data, even a single bit, is detected,
        and the data is rejected.
\end{itemize}
Schemes that provide both are known as \emph{authenticated encryption}. This is the property
modern standards such as TLS~1.3 require.

\textbf{One well-studied building block.} NAPQES is built entirely on \textbf{HMAC-SHA256}, a
standardised and widely deployed cryptographic function (NIST FIPS~198-1 and FIPS~180-4). It does
not depend on the mathematical problems behind RSA or elliptic-curve cryptography, which a large
quantum computer is expected to break. Symmetric schemes keyed with 256-bit secrets, like NAPQES,
are generally considered to remain robust against quantum attacks.

\textbf{How it works, in plain terms:}
\begin{enumerate}
  \item \textbf{Two independent secrets.} Each key has two parts that are generated separately:
        a set of large prime numbers and a 256-bit secret. All security guarantees for the
        current format rest on the 256-bit secret.
  \item \textbf{A unique fingerprint per message.} A message-specific value is derived from the
        message and its context, instead of relying on a random number that must never repeat.
        An observer can tell when exactly the same message is sent twice under the same key and
        context. In exchange, accidental misuse is far less dangerous than in many common
        schemes.
  \item \textbf{Hiding in noise.} Each character becomes a large number, mixed with
        secret-derived decoy values (``noise''), and the whole stream is then masked.
  \item \textbf{Hiding the length.} Messages are padded to fixed size bands, so an observer
        learns at most the band. A stricter profile hides the length entirely.
  \item \textbf{Seal before open.} Every encrypted message carries a cryptographic seal. The seal
        is checked \emph{before} any decryption takes place, so tampered data is rejected
        without revealing anything.
\end{enumerate}

\textbf{Implementations.} Epineon maintains three implementations: Rust (the reference
implementation), Python and C. All three produce identical outputs, and this is checked
automatically against shared test vectors.

\section{Who audited it}
\textbf{ABDK Consulting} is an independent security firm established in 2016. Its team has
conducted more than 300 audits and co-designed widely used cryptographic primitives such as the
Poseidon hash function. The NAPQES reviews were led by \textbf{Dmitry Khovratovich}, a
cryptographer who co-designed the Argon2 password-hashing standard (RFC~9106).

\section{The review rounds}
\begin{tabularx}{\linewidth}{@{}c l >{\raggedright\arraybackslash}p{4.3cm} >{\raggedright\arraybackslash}X@{}}
\toprule
\textbf{Round} & \textbf{Date} & \textbf{What was reviewed} & \textbf{Outcome}\\
\midrule
1 & July 2026 & Specification (research paper), version 1 &
  32 findings: 2 critical, 14 high, 16 moderate or low. Epineon revised the specification in
  response.\\
2 & July 2026 & Specification, version 2 &
  27 of the 32 earlier findings fixed, 3 partly fixed, 2 kept by design. 14 new, mostly
  low-severity findings.\\
3 & August 2026 & Specification, version 3 &
  18 of 19 carried-over findings fixed, 1 partly fixed, none unresolved. The auditors found that
  version~3 \emph{``repairs every soundness defect the previous rounds identified in the security
  analysis''}.\\
4 & Sept.--Oct.\ 2026 & \textbf{Reference software (Rust)}, checked against the specification,
  version 4 &
  37 findings in the initial code. After Epineon's fixes, 25 were fully and 12 partly resolved,
  with none left unresolved and none worse. The follow-up review of the fixed code found
  \textbf{no critical or major issues}. \textbf{Final report v1.0: 27 October 2026.}\\
\bottomrule
\end{tabularx}

After the final review, Epineon addressed every remaining follow-up item. The fixes were
published openly, one change per item. ABDK's final report includes Epineon's response to each
item, and Epineon publishes a point-by-point record of these fixes alongside the report.

\section{What the auditors concluded}
From the final report (ABDK Consulting, v1.0, 27 October 2026):
\pquote{The v8 encryption path is sound, and we found no way to break it from outside.}
\pquote{Overall, the implementation is careful and close to the construction that the paper
describes, and the client's fixes resolved most of the problems we found.}
\pquote{The specification now describes the scheme that the code computes: a re-implementation
from the updated specification reproduced all twenty-two positive v8 test vectors.}

\section{What we ship is what was audited}
\begin{itemize}
  \item \textbf{Specification and software match.} The auditors wrote their own implementation
        from the specification alone. It reproduced the outputs of Epineon's reference software
        byte for byte.
  \item \textbf{Three implementations, one behaviour.} The Python and C implementations are
        tested against the same test vectors as the audited Rust reference and must match it
        exactly.
  \item \textbf{Safe defaults.} A standard build exposes only the audited format (v8). The legacy
        format, kept only for reading archived data, and an experimental streaming mode are
        switched off unless explicitly enabled.
\end{itemize}
\textbf{Further evidence, beyond the audit:}
\begin{itemize}
  \item start-up self-tests modelled on FIPS~140-3;
  \item automated fuzz testing with malformed inputs;
  \item statistical timing tests that check the seal comparison does not leak information;
  \item the NIST SP~800-22 randomness tests applied to encrypted output.
\end{itemize}

\section{Scope and limits --- please read}
\begin{itemize}
  \item \textbf{An audit is a snapshot, not a guarantee.} In ABDK's words, an audit is
        \emph{``not a guarantee of absolute security but rather a snapshot of the system's
        security posture at a specific point in time.''}
  \item \textbf{What was covered.} ABDK reviewed the NAPQES specification in rounds 1 to 3 and
        the Rust reference implementation in round 4. The Python and C implementations were not
        delivered to ABDK; they match the reference through the shared test vectors.
  \item \textbf{What was not covered.} Integration components, such as key exchange, the
        industrial gateway and the protocol adapters, were outside the audit scope.
  \item \textbf{The last round of fixes is self-reported.} ABDK has not verified the fixes
        Epineon made after the final review, nor the latest specification revision (4.2). Both
        are documented publicly so that anyone can check them.
  \item \textbf{No certification is implied.} NAPQES has not been certified under FIPS~140-3 or
        any other scheme. As a recent design, it benefits from continued independent review.
\end{itemize}

\section{Documents and contact}
\begin{tabularx}{\linewidth}{@{}>{\raggedright\arraybackslash}X l@{}}
\toprule
\textbf{Document} & \textbf{Audience}\\
\midrule
ABDK final audit report v1.0 & Technical, full detail\\
Epineon post-audit addendum (fix-by-fix record) & Technical\\
NAPQES specification, revision 4.2 & Technical, cryptographers\\
Source code: \url{https://github.com/Synergeon-epineon/NAPQES} & Engineers\\
\bottomrule
\end{tabularx}

Security questions: \textbf{security@epineon.com}

\vfill
{\footnotesize\color{muted}This summary was prepared by Epineon. It is not an ABDK Consulting
document and has not been reviewed or endorsed by ABDK. Quotations are from ABDK's public report
v1.0. If this summary and that report differ, the report prevails.\par}

\end{document}
